
\documentclass{article}
\usepackage{colortbl}
\usepackage{makecell}
\usepackage{multirow}
\usepackage{supertabular}

\begin{document}

\newcounter{utterance}

\twocolumn

{ \footnotesize  \setcounter{utterance}{1}
\setlength{\tabcolsep}{0pt}
\begin{supertabular}{c@{$\;$}|p{.15\linewidth}@{}p{.15\linewidth}p{.15\linewidth}p{.15\linewidth}p{.15\linewidth}p{.15\linewidth}}

    \# & $\;$A & \multicolumn{4}{c}{Game Master} & $\;\:$B\\
    \hline 

    \theutterance \stepcounter{utterance}  

    & & \multicolumn{4}{p{0.6\linewidth}}{\cellcolor[rgb]{0.9,0.9,0.9}{%
	\makecell[{{p{\linewidth}}}]{% 
	  \tt {\tiny [A$\langle$GM]}  
	 Sono il tuo game master, e stai giocando un gioco collaborativo con la seguente griglia come tavolo da gioco:\\ \tt \\ \tt ```\\ \tt  1234567\\ \tt ╔═══════╗\\ \tt ║◌◌◌◌◌◌◌║ 1\\ \tt ║◌C◌◌◌◌◌║ 2\\ \tt ║◌◌◌◌◌┌─╢ 3\\ \tt ║◌◌◌┌─┤◌║ 4\\ \tt ║◌◌◌│L├─╢ 5\\ \tt ╟───┼─┘◌║ 6\\ \tt ║◌◌◌│◌◌P║ 7\\ \tt ╚═══╧═══╝\\ \tt \\ \tt ```\\ \tt \\ \tt * Il bordo superiore mostra le coordinate x che aumentano verso destra, e il bordo destro le coordinate y che aumentano verso il basso.\\ \tt * I seguenti oggetti sono posizionati casualmente sulla tua griglia: 'C', 'L', 'P'.\\ \tt \\ \tt \\ \tt L'altro giocatore vede una variante del tavolo da gioco, dove gli oggetti sono posizionati in diverse posizioni casuali. Non puoi vedere il tavolo dell'altro giocatore, e loro non possono vedere il tuo.\\ \tt \\ \tt \\ \tt **Obiettivo:**\\ \tt \\ \tt Entrambi i giocatori devono spostare gli oggetti sui rispettivi sfondi in modo che oggetti identici finiscano alle stesse coordinate. Devi comunicare con l'altro giocatore per concordare una configurazione obiettivo comune.\\ \tt \\ \tt \\ \tt **Regole:**\\ \tt \\ \tt * In ogni turno, puoi inviare esattamente uno dei seguenti due comandi:\\ \tt  1. `COMUNICA: <MESSAGE>`: per inviare un messaggio (tutto fino alla prossima interruzione di riga) all'altro giocatore. Lo inoltrerò al tuo partner.\\ \tt 2. `SPOSTA: <OBJECT>, (<X>, <Y>)`: per spostare un oggetto in una nuova posizione, dove <X> è la colonna e <Y> è la riga. Ti farò sapere se la tua mossa è andata a buon fine o no.\\ \tt * Se non rispetti il formato, o invii diversi comandi contemporaneamente, devo penalizzarti.\\ \tt * Se entrambi i giocatori accumulano più di 8 penalità, perdete entrambi il gioco.\\ \tt * È vitale che comunichi con l'altro giocatore riguardo al tuo stato obiettivo! L'*unico* modo in cui puoi trasmettere la tua strategia all'altro giocatore è usando il comando `COMUNICA: <MESSAGE>`!\\ \tt \\ \tt \\ \tt **Spostamento Oggetti**\\ \tt \\ \tt * Puoi spostare oggetti solo in celle entro i limiti della griglia. La cella di destinazione deve essere vuota, cioè deve contenere solo il simbolo '◌'.\\ \tt * Se provi a spostare un oggetto in un posto che non è vuoto, o provi a spostarlo fuori dalla griglia, devo penalizzarti. Hai un altro tentativo.\\ \tt * Prima di fare una mossa, controlla due volte che il posto di destinazione sia vuoto, e non contenga alcuna lettera, cornice, o linea!\\ \tt \\ \tt \\ \tt **Fine del Gioco**\\ \tt \\ \tt Se pensi di aver raggiunto l'obiettivo di allineare tutti gli oggetti, puoi chiedere all'altro giocatore di finire il gioco inviando COMUNICA: finito?. Se l'altro giocatore ti chiede di finire il gioco, e rispondi con COMUNICA: finito!, the game will end.\\ \tt \\ \tt Entrambi i giocatori vincono se il gioco finisce entro 12 round, dove un round è definito come due giocatori che inviano ciascuno un comando valido.\\ \tt \\ \tt **Punteggio:**\\ \tt \\ \tt Più vicini sono gli oggetti identici in entrambi i tavoli da gioco, più punti ottieni. Le penalità riducono i tuoi punti. Riesci a battere il record?\\ \tt \\ \tt \\ \tt Per favore invia un messaggio all'altro giocatore per iniziare il gioco! 
	  } 
	   } 
	   } 
	 & \\ 
 

    \theutterance \stepcounter{utterance}  

    & \multicolumn{4}{p{0.6\linewidth}}{\cellcolor[rgb]{0.8,1,0.9}{%
	\makecell[{{p{\linewidth}}}]{% 
	  \tt {\tiny [A$\rangle$GM]}  
	 COMUNICA: Ciao! Posizioni sul mio tavolo: C(2,2), L(5,5), P(7,7). Propongo come obiettivo comune: C→(2,2), L→(5,5), P→(7,7). Queste celle sono libere da pareti qui. Confermi? Se ok, muovi i tuoi pezzi lì; io resto fermo. Quando sei allineato, scrivi “finito?”. Se non va, proponi tu una terna alternativa. 
	  } 
	   } 
	   } 
	 & & \\ 
 

    \theutterance \stepcounter{utterance}  

    & & & \multicolumn{2}{p{0.3\linewidth}}{\cellcolor[rgb]{0.95,0.95,0.95}{%
	\makecell[{{p{\linewidth}}}]{% 
	  \tt {\tiny [GM$|$GM]}  
	 COMUNICA: Ciao! Posizioni sul mio tavolo: C(2,2), L(5,5), P(7,7). Propongo come obiettivo comune: C→(2,2), L→(5,5), P→(7,7). Queste celle sono libere da pareti qui. Confermi? Se ok, muovi i tuoi pezzi lì; io resto fermo. Quando sei allineato, scrivi “finito?”. Se non va, proponi tu una terna alternativa. 
	  } 
	   } 
	   } 
	 & & \\ 
 

    \theutterance \stepcounter{utterance}  

    & & \multicolumn{4}{p{0.6\linewidth}}{\cellcolor[rgb]{0.9,0.9,0.9}{%
	\makecell[{{p{\linewidth}}}]{% 
	  \tt {\tiny [GM$\rangle$B]}  
	 Sono il tuo game master, e stai giocando un gioco collaborativo con la seguente griglia come tavolo da gioco:\\ \tt \\ \tt ```\\ \tt  1234567\\ \tt ╔═══════╗\\ \tt ║◌◌◌◌◌◌◌║ 1\\ \tt ║◌◌C◌◌◌◌║ 2\\ \tt ║◌◌◌◌◌┌─╢ 3\\ \tt ║◌◌◌┌─┤◌║ 4\\ \tt ║◌P◌│◌├─╢ 5\\ \tt ╟───┼─┘◌║ 6\\ \tt ║L◌◌│◌◌◌║ 7\\ \tt ╚═══╧═══╝\\ \tt \\ \tt ```\\ \tt \\ \tt * Il bordo superiore mostra le coordinate x che aumentano verso destra, e il bordo destro le coordinate y che aumentano verso il basso.\\ \tt * I seguenti oggetti sono posizionati casualmente sulla tua griglia: 'C', 'L', 'P'.\\ \tt \\ \tt \\ \tt L'altro giocatore vede una variante del tavolo da gioco, dove gli oggetti sono posizionati in diverse posizioni casuali. Non puoi vedere il tavolo dell'altro giocatore, e loro non possono vedere il tuo.\\ \tt \\ \tt \\ \tt **Obiettivo:**\\ \tt \\ \tt Entrambi i giocatori devono spostare gli oggetti sui rispettivi sfondi in modo che oggetti identici finiscano alle stesse coordinate. Devi comunicare con l'altro giocatore per concordare una configurazione obiettivo comune.\\ \tt \\ \tt \\ \tt **Regole:**\\ \tt \\ \tt * In ogni turno, puoi inviare esattamente uno dei seguenti due comandi:\\ \tt  1. `COMUNICA: <MESSAGE>`: per inviare un messaggio (tutto fino alla prossima interruzione di riga) all'altro giocatore. Lo inoltrerò al tuo partner.\\ \tt 2. `SPOSTA: <OBJECT>, (<X>, <Y>)`: per spostare un oggetto in una nuova posizione, dove <X> è la colonna e <Y> è la riga. Ti farò sapere se la tua mossa è andata a buon fine o no.\\ \tt * Se non rispetti il formato, o invii diversi comandi contemporaneamente, devo penalizzarti.\\ \tt * Se entrambi i giocatori accumulano più di 8 penalità, perdete entrambi il gioco.\\ \tt * È vitale che comunichi con l'altro giocatore riguardo al tuo stato obiettivo! L'*unico* modo in cui puoi trasmettere la tua strategia all'altro giocatore è usando il comando `COMUNICA: <MESSAGE>`!\\ \tt \\ \tt \\ \tt **Spostamento Oggetti**\\ \tt \\ \tt * Puoi spostare oggetti solo in celle entro i limiti della griglia. La cella di destinazione deve essere vuota, cioè deve contenere solo il simbolo '◌'.\\ \tt * Se provi a spostare un oggetto in un posto che non è vuoto, o provi a spostarlo fuori dalla griglia, devo penalizzarti. Hai un altro tentativo.\\ \tt * Prima di fare una mossa, controlla due volte che il posto di destinazione sia vuoto, e non contenga alcuna lettera, cornice, o linea!\\ \tt \\ \tt \\ \tt **Fine del Gioco**\\ \tt \\ \tt Se pensi di aver raggiunto l'obiettivo di allineare tutti gli oggetti, puoi chiedere all'altro giocatore di finire il gioco inviando COMUNICA: finito?. Se l'altro giocatore ti chiede di finire il gioco, e rispondi con COMUNICA: finito!, the game will end.\\ \tt \\ \tt Entrambi i giocatori vincono se il gioco finisce entro 12 round, dove un round è definito come due giocatori che inviano ciascuno un comando valido.\\ \tt \\ \tt **Punteggio:**\\ \tt \\ \tt Più vicini sono gli oggetti identici in entrambi i tavoli da gioco, più punti ottieni. Le penalità riducono i tuoi punti. Riesci a battere il record?\\ \tt \\ \tt L'altro giocatore ha iniziato il gioco inviando questo messaggio:\\ \tt \\ \tt "Ciao! Posizioni sul mio tavolo: C(2,2), L(5,5), P(7,7). Propongo come obiettivo comune: C→(2,2), L→(5,5), P→(7,7). Queste celle sono libere da pareti qui. Confermi? Se ok, muovi i tuoi pezzi lì; io resto fermo. Quando sei allineato, scrivi “finito?”. Se non va, proponi tu una terna alternativa."\\ \tt \\ \tt Qual è il tuo primo comando? 
	  } 
	   } 
	   } 
	 & \\ 
 

    \theutterance \stepcounter{utterance}  

    & & & \multicolumn{4}{p{0.6\linewidth}}{\cellcolor[rgb]{1,0.85,0.72}{%
	\makecell[{{p{\linewidth}}}]{% 
	  \tt {\tiny [GM$\langle$B]}  
	 COMUNICA: Confermo. Obiettivo comune: C(2,2), L(5,5), P(7,7). Quelle celle sono libere sulla mia griglia. Muovo subito i miei pezzi lì (prima C, poi L, poi P). Resta fermo; quando avrò finito ti mando “finito?”. 
	  } 
	   } 
	   } 
	  \\ 
 

    \theutterance \stepcounter{utterance}  

    & & & \multicolumn{2}{p{0.3\linewidth}}{\cellcolor[rgb]{0.95,0.95,0.95}{%
	\makecell[{{p{\linewidth}}}]{% 
	  \tt {\tiny [GM$|$GM]}  
	 COMUNICA: Confermo. Obiettivo comune: C(2,2), L(5,5), P(7,7). Quelle celle sono libere sulla mia griglia. Muovo subito i miei pezzi lì (prima C, poi L, poi P). Resta fermo; quando avrò finito ti mando “finito?”. 
	  } 
	   } 
	   } 
	 & & \\ 
 

    \theutterance \stepcounter{utterance}  

    & & \multicolumn{4}{p{0.6\linewidth}}{\cellcolor[rgb]{0.9,0.9,0.9}{%
	\makecell[{{p{\linewidth}}}]{% 
	  \tt {\tiny [A$\langle$GM]}  
	 Il tuo messaggio è stato inoltrato all'altro giocatore.\\ \tt \\ \tt \\ \tt Stai attualmente giocando il round 1 di massimo 12.\\ \tt Avete accumulato collettivamente 0 di 8 penalità.\\ \tt L'altro giocatore ha inviato questo messaggio:\\ \tt \\ \tt \\ \tt "Confermo. Obiettivo comune: C(2,2), L(5,5), P(7,7). Quelle celle sono libere sulla mia griglia. Muovo subito i miei pezzi lì (prima C, poi L, poi P). Resta fermo; quando avrò finito ti mando “finito?”."\\ \tt \\ \tt \\ \tt Qual è il tuo prossimo comando? 
	  } 
	   } 
	   } 
	 & \\ 
 

    \theutterance \stepcounter{utterance}  

    & \multicolumn{4}{p{0.6\linewidth}}{\cellcolor[rgb]{0.8,1,0.9}{%
	\makecell[{{p{\linewidth}}}]{% 
	  \tt {\tiny [A$\rangle$GM]}  
	 COMUNICA: Perfetto. Confermo: non muovo nulla. Quando hai allineato C(2,2), L(5,5), P(7,7), scrivi "finito?" e ti rispondo "finito!". 
	  } 
	   } 
	   } 
	 & & \\ 
 

    \theutterance \stepcounter{utterance}  

    & & & \multicolumn{2}{p{0.3\linewidth}}{\cellcolor[rgb]{0.95,0.95,0.95}{%
	\makecell[{{p{\linewidth}}}]{% 
	  \tt {\tiny [GM$|$GM]}  
	 COMUNICA: Perfetto. Confermo: non muovo nulla. Quando hai allineato C(2,2), L(5,5), P(7,7), scrivi "finito?" e ti rispondo "finito!". 
	  } 
	   } 
	   } 
	 & & \\ 
 

    \theutterance \stepcounter{utterance}  

    & & & \multicolumn{2}{p{0.3\linewidth}}{\cellcolor[rgb]{0.95,0.95,0.95}{%
	\makecell[{{p{\linewidth}}}]{% 
	  \tt {\tiny [GM$|$GM]}  
	 true 
	  } 
	   } 
	   } 
	 & & \\ 
 

    \theutterance \stepcounter{utterance}  

    & & & \multicolumn{2}{p{0.3\linewidth}}{\cellcolor[rgb]{0.95,0.95,0.95}{%
	\makecell[{{p{\linewidth}}}]{% 
	  \tt {\tiny [GM$|$GM]}  
	 * success: True\\ \tt * lose: False\\ \tt * aborted: False\\ \tt -------\\ \tt * Shifts: 0.00\\ \tt * Max Shifts: 4.00\\ \tt * Min Shifts: 2.00\\ \tt * End Distance Sum: 10.86\\ \tt * Init Distance Sum: 10.86\\ \tt * Expected Distance Sum: 12.57\\ \tt * Penalties: 0.00\\ \tt * Max Penalties: 8.00\\ \tt * Rounds: 1.00\\ \tt * Max Rounds: 12.00\\ \tt * Object Count: 3.00\\ \tt  
	  } 
	   } 
	   } 
	 & & \\ 
 

    \theutterance \stepcounter{utterance}  

    & & & \multicolumn{2}{p{0.3\linewidth}}{\cellcolor[rgb]{0.95,0.95,0.95}{%
	\makecell[{{p{\linewidth}}}]{% 
	  \tt {\tiny [GM$|$GM]}  
	 * Main Score: 8.38\\ \tt \\ \tt -------\\ \tt * Distance Score: 0.07\\ \tt * Coverage Score: 0.06\\ \tt * Penalty Score: 1.29\\ \tt * Alternative Penalty Score: 1.12\\ \tt * Alternative Main Score: 7.67\\ \tt \\ \tt -------\\ \tt * Shifts: 0.00\\ \tt * Max Shifts: 4.00\\ \tt * Min Shifts: 2.00\\ \tt * End Distance Sum: 10.86\\ \tt * Init Distance Sum: 10.86\\ \tt * Expected Distance Sum: 12.57\\ \tt * Penalties: 0.00\\ \tt * Max Penalties: 8.00\\ \tt * Rounds: 1.00\\ \tt * Max Rounds: 12.00\\ \tt * Object Count: 3.00\\ \tt  
	  } 
	   } 
	   } 
	 & & \\ 
 

\end{supertabular}
}

\end{document}
